\documentclass[aspectratio=169,11pt]{beamer}
% Shared Beamer style for practice contest solution walkthroughs.
\usepackage[utf8]{inputenc}
\usepackage[T1]{fontenc}
\usepackage{lmodern}
\usepackage{amsmath,amssymb}
\usepackage{xcolor}
\usepackage{hyperref}
\usepackage{listings}

\definecolor{CPNavy}{HTML}{172033}
\definecolor{CPBlue}{HTML}{2563EB}
\definecolor{CPGreen}{HTML}{16A34A}
\definecolor{CPOrange}{HTML}{F97316}
\definecolor{CPGray}{HTML}{64748B}
\definecolor{CPLight}{HTML}{F8FAFC}
\definecolor{CPCodeBg}{HTML}{F1F5F9}
\definecolor{CPCodeRule}{HTML}{CBD5E1}

\usetheme{default}
\usefonttheme{professionalfonts}
\setbeamertemplate{navigation symbols}{}
\setbeamersize{text margin left=0.65cm,text margin right=0.65cm}
\setbeamertemplate{itemize item}{\color{CPBlue}$\blacktriangleright$}
\setbeamertemplate{itemize subitem}{\color{CPGray}$\triangleright$}

\setbeamercolor{normal text}{fg=CPNavy,bg=white}
\setbeamercolor{frametitle}{fg=white,bg=CPNavy}
\setbeamercolor{title}{fg=white,bg=CPNavy}
\setbeamercolor{subtitle}{fg=CPLight,bg=CPNavy}
\hypersetup{colorlinks=true,urlcolor=CPBlue}

\setbeamertemplate{frametitle}{%
  \nointerlineskip%
  \begin{beamercolorbox}[wd=\paperwidth,ht=0.88cm,dp=0.22cm,leftskip=0.55cm,rightskip=0.35cm]{frametitle}
    \usebeamerfont{frametitle}\insertframetitle\hfill\small\insertframenumber/\inserttotalframenumber
  \end{beamercolorbox}%
}

\setbeamertemplate{footline}{%
  \leavevmode%
  \hbox{%
    \begin{beamercolorbox}[wd=.58\paperwidth,ht=2.4ex,dp=1ex,leftskip=.5cm]{author in head/foot}%
      \color{CPGray}\insertshorttitle
    \end{beamercolorbox}%
    \begin{beamercolorbox}[wd=.42\paperwidth,ht=2.4ex,dp=1ex,rightskip=.5cm plus1fil]{date in head/foot}%
      \hfill\color{CPGray}\insertshortauthor
    \end{beamercolorbox}%
  }%
}

\setbeamertemplate{title page}{%
  \vspace*{1.25cm}
  \begin{beamercolorbox}[wd=\paperwidth,sep=0.65cm,leftskip=0.15cm,rightskip=0.15cm]{title}
    {\color{white}\Huge\bfseries\inserttitle\par}
    \vspace{0.35cm}
    {\color{CPLight}\Large\insertsubtitle\par}
    \vspace{0.6cm}
    {\color{CPLight}\large\insertauthor\par}
  \end{beamercolorbox}
  \vfill
}

\newcommand{\sectiontag}[1]{\textbf{\color{CPBlue}#1}}
\newenvironment{tightitemize}{\begin{itemize}\small\setlength{\itemsep}{0.18cm}}{\end{itemize}}
\lstdefinestyle{cpcode}{
  language=Python,
  basicstyle=\ttfamily\tiny,
  keywordstyle=\color{CPBlue}\bfseries,
  commentstyle=\color{CPGray}\itshape,
  stringstyle=\color{CPGreen},
  backgroundcolor=\color{CPCodeBg},
  frame=single,
  framerule=0.25pt,
  rulecolor=\color{CPCodeRule},
  showstringspaces=false,
  columns=fullflexible,
  keepspaces=true,
  breaklines=true,
  breakatwhitespace=false,
  tabsize=4,
  xleftmargin=0.08cm,
  xrightmargin=0.08cm,
  aboveskip=0.08cm,
  belowskip=0.08cm
}


\title[convexpolygonarea]{Convex Polygon Area}
\subtitle{Day 5 solution walkthrough}
\author[Solution Deck]{Competitive Programming Bootcamp}
\date{}

\begin{document}

\begin{frame}[plain]
  \titlepage
\end{frame}

% BEGIN GENERATED STATEMENT INTERPRETATION
\begin{frame}{Interpret the Word Problem}
\small
\sectiontag{Statement excerpts:}
\begin{tightitemize}
  \item \textit{``calculate the area of convex polygons''}
  \item \textit{``the last vertex is connected to the first one''}
\end{tightitemize}

\vspace{0.18cm}
\sectiontag{Programming interpretation:}
\begin{tightitemize}
  \item The input already gives vertices in polygon order.
  \item The shoelace formula sums each directed edge contribution.
  \item Twice the area is an integer, so the final answer ends in .0 or .5.
\end{tightitemize}
\end{frame}
% END GENERATED STATEMENT INTERPRETATION







\begin{frame}{Problem Model}
\small
\sectiontag{Textbook connection:} CP4 7.3.3 Area of a Polygon

\vspace{0.25cm}
\sectiontag{Core technique:} Shoelace formula

\vspace{0.25cm}
\begin{tightitemize}
  \item Each test case gives a polygon as ordered vertices.
  \item The polygon is convex, but the shoelace formula works for any simple ordered polygon.
  \item The answer is the area printed with .0 or .5.
\end{tightitemize}
\end{frame}

\begin{frame}{How to Recognize the Approach}
\begin{tightitemize}
  \item Ordered polygon vertices suggest summing edge contributions.
  \item Connecting the last vertex back to the first is required.
  \item Integer coordinates make twice the area an integer.
\end{tightitemize}
\end{frame}

\begin{frame}{Algorithm}
\begin{tightitemize}
  \item Parse the input line into n coordinate pairs.
  \item Initialize area2 = 0.
  \item For each vertex i, let j = (i + 1) mod n.
  \item Add x\_i*y\_j - y\_i*x\_j to area2.
  \item Take abs(area2), divide by 2, and print .0 or .5.
\end{tightitemize}
\end{frame}

\begin{frame}{Walkthrough of the Key State}
\begin{tightitemize}
  \item area2 is twice the signed area.
  \item The modulo index closes the polygon without a special last-edge case.
  \item The sign depends on clockwise versus counterclockwise order, so take absolute value.
\end{tightitemize}
\end{frame}

\begin{frame}{Why the Algorithm Is Correct}
\begin{tightitemize}
  \item The shoelace formula sums the signed area contribution of each directed polygon edge.
  \item Closing the polygon includes every boundary edge exactly once.
  \item Taking the absolute value removes orientation, yielding the geometric area.
\end{tightitemize}
\end{frame}

\begin{frame}{Complexity and Implementation Notes}
\small
\sectiontag{Complexity:} O(N) time and O(N) memory per polygon.

\vspace{0.3cm}
\begin{tightitemize}
  \item You can compute while parsing to reduce memory, but storing points keeps the code clear.
  \item The final area is always an integer or half-integer for integer coordinates.
\end{tightitemize}
\end{frame}



% BEGIN GENERATED CODE WALKTHROUGH
\begin{frame}[fragile]{Code: Parse One Polygon}
\scriptsize
\sectiontag{Source lines:} \texttt{convexpolygonearea.py:41--65}

\vspace{0.08cm}
\begin{columns}[T,totalwidth=\textwidth]
\begin{column}{0.58\textwidth}
\begin{lstlisting}[style=cpcode]
import sys

def main():
    input = sys.stdin.buffer.readline

    test_cases = int(input())

    for _ in range(test_cases):

        values = list(map(int, input().split()))

        n = values[0]

        points = []

        for i in range(n):
            x = values[1 + 2 * i]
            y = values[2 + 2 * i]

            points.append((x, y))
\end{lstlisting}
\end{column}
\begin{column}{0.38\textwidth}
\sectiontag{What this chunk does}
\begin{tightitemize}
  \item Each polygon is described on one line as n followed by coordinate pairs.
  \item The loop converts the flat coordinate list into (x, y) points.
  \item The points stay in the counter-clockwise order supplied by the input.
\end{tightitemize}
\end{column}
\end{columns}
\end{frame}

\begin{frame}[fragile]{Code: Apply Shoelace Formula}
\scriptsize
\sectiontag{Source lines:} \texttt{convexpolygonearea.py:67--84}

\vspace{0.08cm}
\begin{columns}[T,totalwidth=\textwidth]
\begin{column}{0.58\textwidth}
\begin{lstlisting}[style=cpcode]
area2 = 0

for i in range(n):
    x1, y1 = points[i]

    x2, y2 = points[(i + 1) % n]

    area2 += x1 * y2 - y1 * x2

area2 = abs(area2)
\end{lstlisting}
\end{column}
\begin{column}{0.38\textwidth}
\sectiontag{What this chunk does}
\begin{tightitemize}
  \item area2 stores twice the signed area.
  \item Modulo indexing closes the edge from the last vertex back to the first.
  \item The absolute value removes dependence on orientation.
\end{tightitemize}
\end{column}
\end{columns}
\end{frame}

\begin{frame}[fragile]{Code: Print Half-Integer Area}
\scriptsize
\sectiontag{Source lines:} \texttt{convexpolygonearea.py:86--94}

\vspace{0.08cm}
\begin{columns}[T,totalwidth=\textwidth]
\begin{column}{0.58\textwidth}
\begin{lstlisting}[style=cpcode]
        if area2 % 2 == 0:
            print(f"{area2 // 2}.0")
        else:
            print(f"{area2 // 2}.5")

main()
\end{lstlisting}
\end{column}
\begin{column}{0.38\textwidth}
\sectiontag{What this chunk does}
\begin{tightitemize}
  \item Even area2 values produce an integer area ending in .0.
  \item Odd area2 values produce a half-integer area ending in .5.
  \item main() runs the full parser and area calculation.
\end{tightitemize}
\end{column}
\end{columns}
\end{frame}
% END GENERATED CODE WALKTHROUGH

\begin{frame}{Source}
\small
\sectiontag{Solution file:} \texttt{practice\_contests/day\_5/convexpolygonearea.py}

\vspace{0.3cm}
\sectiontag{Problem:} \url{https://open.kattis.com/problems/convexpolygonarea}

\vspace{0.45cm}
Use this deck with the source code open beside it. The slides explain the reasoning; the code shows the exact input handling and output formatting.
\end{frame}

\end{document}
